\scanpage{264}
\mspage{214}

Pour le dire autrement, on a défini \uline{deux} notions de \inserted{$W$-}types d'homotopie relatifs sur $S$, l'une de façon à donner un tel type d'homotopie \inserted{$\mathrm{ho}^{!}_{W,S}(X)$} pour \uline{tout} objet de $M_{/S}$, mais il n'y a \inserted{a priori} une bonne notion de chgt de base que pour les chgt. de base $S'\unsure{\to|\text{---}}S$ qui sont $W$-fibrants; l'autre, qui correspond plus directement à l'intuition d'un $W$-type d'homotopie loc. constant sur $S$, n'associe un tel type $\mathrm{ho}^{\unsure{\mathrm{lc}|\mathrm{le}}}_{W,S}(X)$ qu'aux $X$ de $M_{/S}$ qui sont $W$-fibrés sur $S$, en revanche \unsure{cette} notion admet une bonne notion d'image inverse et de ``fibre $W$-homotopique'' pour des chgt. de base quelconques.

Une ``bonne'' théorie des fibres $W$-homotopiques doit réunir en une seule les deux notions précédentes, de sorte que les avantages des deux s'ajoutent, et les désavantages s'éliminent ! Elle s'obtient à l'aide d'un axiome draconien sur la situation (\unsure{alors|donc} que jusqu'à maintenant on n'a rien utilisé sur $W\subset\mathrm{Fl}(M)$, si ce n'est \deleted{\unsure{tout}} au plus la saturation faible) \unsure{,|;} en plus de l'existence des produits fibrés)~:
% archive sheet number: 264
